\section{Data}
\label{sec:data}

We use the JPX Tokyo Stock Exchange daily price panel, covering
\num{2000} securities. The training file spans 2017-01-04 to 2021-12-03 and the
supplemental file continues to 2022-06-24; after feature construction the
usable panel holds \num{2483470} security-days over \num{1296} trading dates
from 2017-02-02. The \num{2000} securities are the competition's primary
universe, a membership list fixed by the organiser when the data were released
at the end of the training period and applied to the whole window. It therefore
carries the hindsight that Section~\ref{sec:intro} criticises: a name is in the
panel for 2017 because it was listed and liquid in late 2021. The bias is common
to every condition and to the published result we re-score, so it enters no
contrast between them, and no absolute level of performance on this panel is
claimed.

\subsection{Split adjustment}
\label{sec:data:adjustment}
The panel records a per-day \texttt{AdjustmentFactor}. A factor of $0.5$ on day
$t$ indicates that prices up to and including $t$ must be halved to be
comparable with prices after $t$, so the cumulative factor at date $t$ is the
product of all factors from $t$ onwards:
\begin{equation}
  \tilde{P}_t = P_t \prod_{s \ge t} a_s .
\end{equation}
Publicly circulated pipelines for this dataset instead accumulate forward,
$\prod_{s \le t} a_s$. Table~\ref{tab:adjustment} shows the consequence over the
\num{610} securities with at least one non-unit factor.

\begin{table}[t]
\centering
\caption{Artefactual daily moves under three adjustment conventions.}
\label{tab:adjustment}
\begin{tabular}{lrrr}
\toprule
Adjustment & $|r|>30\%$ & $|r|>50\%$ & $\max |r|$ \\
\midrule
None (raw closes)            & 718  & 565  & 20.09 \\
Forward cumulative product   & 1418 & 1091 & 23.43 \\
Backward cumulative product  & \textbf{16} & \textbf{3} & \textbf{1.12} \\
\bottomrule
\end{tabular}
\end{table}

\subsection{Features}
Twenty-one per-security features are computed from the adjusted close $\tilde P_t$,
the daily return $r_t=\tilde P_t/\tilde P_{t-1}-1$, the raw volume $V_t$ and the
dollar volume $D_t=P_tV_t$ (Table~\ref{tab:features}). Every window looks
strictly backwards, so a feature dated $t$ uses closes up to and including $t$
only. The set is deliberately generic (momentum, realised volatility,
standardised returns, liquidity, range position, drawdown and RSI) so that
no model has access to information the others lack.

\begin{table}[t]
  \centering
  \caption{The 21 input features. All rolling statistics use a window of
  exactly $n$ trading days and are undefined until $n$ observations exist.}
  \label{tab:features}
  \small
  \begin{tabular}{llc}
    \toprule
    Group & Definition & $n$ \\
    \midrule
    Momentum ($\mathrm{mom}_n$) & $\tilde P_t/\tilde P_{t-n}-1$ & 1, 2, 5, 10, 20 \\
    Volatility ($\mathrm{vol}_n$) & rolling standard deviation of $r$ & 5, 10, 20 \\
    Standardised return ($\mathrm{ret\_z}_n$) & $(r_t-\bar r_n)/\sigma_n(r)$ & 5, 10, 20 \\
    Log volume & $\log(1+V_t)$ & -- \\
    Log dollar volume & $\log(1+D_t)$ & -- \\
    Dollar-volume $z$-score ($\mathrm{dv\_z}_n$) & rolling $z$-score of $\log(1+D_t)$ & 5, 10, 20 \\
    Range position ($\mathrm{pos}_n$) & $(\tilde P_t-\min_n)/(\max_n-\min_n)$ & 10, 20 \\
    Drawdown ($\mathrm{dd}_{20}$) & $\tilde P_t/\max_{20}-1$ & 20 \\
    RSI ($\mathrm{rsi}_n$) & simple-moving-average RSI, scaled to $[-0.5,0.5]$ & 7, 14 \\
    \bottomrule
  \end{tabular}
\end{table}

\subsection{Cross-sectional normalisation}
Each feature is winsorised at the 1st and 99th percentile and standardised
within each trading date. Both operations use only same-day information and
therefore introduce no look-ahead.

\subsection{Second market}
\label{sec:data-us}
To test whether the result is a property of one exchange, the identical
pipeline is run on United States common stocks over the same calendar window.
Candidates are the \num{4826} NYSE, NYSE American and Nasdaq common stocks in
the exchange symbol directories after removing ETFs, warrants, units, rights,
preferred shares, notes, depositary receipts, trusts and acquisition shells by
name; \num{3935} of them traded in the window. The universe is the \num{2000}
names with the highest median dollar volume in 2017, the first calendar year,
among those with at least \num{200} trading days and a median close of at least
\$2 that year, so selection uses no information after the first training
window. Split-adjusted closes and volumes come from a public daily-price source and
are adjusted by the convention of Section~\ref{sec:data:adjustment}; the
target, the 21 features and the
cross-sectional normalisation are computed by the same functions. The panel
holds \num{2796919} security-days over \num{1399} trading dates from
2016-11-30 to 2022-06-22. Because the US calendar has more trading days, the
development period yields five walk-forward folds (\num{625} out-of-sample
days) rather than four; the holdout window is the same. The symbol directory
is a current snapshot, so names delisted before it are absent; every condition
is trained and scored on this identical panel, so the survivorship bias is
common to all conditions and does not enter the contrasts between them,
and absolute performance levels on this market are not claimed.
